\begingroup
\csname tagsleft@true\endcsname
% Printed 192; principal PDF 199; Google PDF 203.
\par\clearpage
\noindent\makebox[0pt][l]{\fontsize{13}{16}\selectfont 192}\hfill{\small\textbf{POSTHUMOUS PAPERS.}}\hfill\null\par\medskip

From this one finally derives (in our notation)
\[
\begin{aligned}
\log\operatorname{tang}(45^\circ+\tfrac12\omega)&=\log\operatorname{tang}(45^\circ+\tfrac12\psi)\\
&\phantom{={}}-\left(4ff+4f^4-\frac43f^6-132f^8\ldots\right)\sin\psi\\
&\phantom{={}}-\left(\frac83f^4+\frac{52}3f^6+104f^8\ldots\right)\sin3\psi\\
&\phantom{={}}-\left(\frac{88}{15}f^6+\frac{232}3f^8\ldots\right)\sin5\psi\\
&\phantom{={}}-\left(\frac{400}{21}f^8\ldots\right)\sin7\psi.
\end{aligned}
\]
[Now, for the flattening \(\frac{1}{302{,}7827\ldots}\), p.\ 183 gives:]
{\small\setlength{\arraycolsep}{1.5pt}
\[
\begin{array}{rr@{\quad}rr@{\quad}rr@{\quad}rr}
4ff={}&10495{,}452\,5992\,\mathrm m\\
4f^4={}&4{,}325\,7581\phantom{\,\mathrm m}
 & \dfrac83 f^4={}&2{,}883\,8387\,\mathrm m\\[6pt]
-\dfrac43 f^6={}&-0{,}000\,5943\phantom{\,\mathrm m}
 & \dfrac{52}3 f^6={}&0{,}007\,7258\phantom{\,\mathrm m}
 & \dfrac{88}{15}f^6={}&0{,}002\,6149\,\mathrm m\\[6pt]
-132f^8={}&-0{,}000\,0242\phantom{\,\mathrm m}
 & 104f^8={}&0{,}000\,0191\phantom{\,\mathrm m}
 & \dfrac{232}3 f^8={}&0{,}000\,0142\phantom{\,\mathrm m}
 & \dfrac{400}{21}f^8={}&0{,}000\,0035\,\mathrm m\\
 &\overline{10499{,}777\,7388}\,\mathrm m
 & &\overline{2{,}891\,5836}\,\mathrm m
 & &\overline{0{,}002\,6291}\,\mathrm m
 & &\overline{0{,}000\,0035}\,\mathrm m\\[.6ex]
\multicolumn{8}{@{}l}{\text{[or] in arcseconds [(}1\,\mathrm m=0^{\prime\prime}{,}0324\text{)]}}\\[.6ex]
 &340^{\prime\prime}{,}192\,7987\phantom{\,\mathrm m}
 & &0^{\prime\prime}{,}093\,6873\phantom{\,\mathrm m}
 & &0^{\prime\prime}{,}000\,0852\phantom{\,\mathrm m}
 & &0^{\prime\prime}{,}000\,0001.\phantom{\mathrm m}
\end{array}
\]
\par}
[Since] we have
\[
[\operatorname{tang}(45^\circ+\tfrac12(x+iy))=
\frac{\cos x+\sin iy}{\cos iy-\sin x}=
\frac{\cos iy+\sin x}{\cos x-\sin iy},
\]
\noindent it follows]
\[
\log\operatorname{tang}(45^\circ+\tfrac12(x+iy))
=\log\surd\frac{\cos iy+\sin x}{\cos iy-\sin x}
+i\operatorname{arc}\operatorname{tang}\frac{\sin iy}{i\cos x}.
\]
If, then, \(\Omega\) corresponds to the latitude of the place whose coordinates are \(x,y\),
and if \(\lambda\) is its longitude, we have
\[
\begin{aligned}
\mathopen{[}\log\operatorname{tang}(45^\circ+\tfrac12\Omega)+i\lambda
  &=\log\operatorname{tang}(45^\circ+\tfrac12(x+iy))\\
  &\quad-D\sin(x+iy)-E\sin3(x+iy)-F\sin5(x+iy)\ldots,
\end{aligned}
\]
\noindent where \(D,E,F,\ldots\) are the coefficients of the series for
\(\log\operatorname{tang}(45^\circ+\tfrac12\omega)\). If \(D,E,F,\ldots\)
are given in arcseconds, then with \(\rho=\tfrac{1}{206265}\) and for
Briggsian logarithms of modulus \(k\), we have:
\[
\begin{aligned}
\lambda&=\frac1\rho\operatorname{arc}\operatorname{tang}\frac{\sin iy}{i\cos x}
-D\cos x\frac{\sin iy}{i}-E\cos3x\frac{\sin3iy}{i}-\ldots\\
\log\operatorname{tang}(45^\circ+\tfrac12\Omega)
&=\log\surd\frac{\cos iy+\sin x}{\cos iy-\sin x}
-k\rho(D\sin x\cos iy+E\sin3x\cos3iy+\ldots).
\end{aligned}
\]

% Printed 193; principal PDF 200; Google PDF 204.
\par\clearpage
\noindent\null\hfill{\small\textbf{CONFORMAL MAPPING OF THE SPHEROID IN THE PLANE.}}\hfill\makebox[0pt][r]{\fontsize{13}{16}\selectfont 193}\par\medskip

With the stated values of \(D,E,F,\ldots\), we therefore obtain:]
\[
\begin{aligned}
\lambda=\frac1\rho\operatorname{arc}\operatorname{tang}\frac{\sin iy}{i\cos x}
&-340^{\prime\prime}{,}192\,7987\cos x\frac{\sin iy}{i}\\
&-\phantom{34}0^{\prime\prime}{,}093\,6873\cos3x\frac{\sin3iy}{i}\\
&-\phantom{34}0^{\prime\prime}{,}000\,0852\cos5x\frac{\sin5iy}{i}\\
&-\phantom{34}0^{\prime\prime}{,}000\,0001\cos7x\frac{\sin7iy}{i}\\
&\quad\text{etc.}
\end{aligned}
\]
\[
\begin{aligned}
\log\operatorname{tang}(45^\circ+\tfrac12\Omega)
&=\log\surd\frac{\cos iy+\sin x}{\cos iy-\sin x}\\
&\phantom{={}}-0{,}000\,716\,2824\sin x\cos iy\\
&\phantom{={}}-0{,}000\,000\,1973\sin3x\cos3iy\\
&\phantom{={}}-0{,}000\,000\,0002\sin5x\cos5iy\\
&\qquad\qquad\text{etc.}
\end{aligned}
\]
[Using the table on p.\ 184, we then obtain for \(\Omega\) the corresponding \(\Phi\).]

\begin{center}[32.]\end{center}

[The representation of the surface of the spheroid in the plane is] expressed by the following two formulas:
\begin{equation*}
\operatorname{cotang}\tfrac12 p\left(\frac{1-e\cos p}{1+e\cos p}\right)^{\frac12 e}
\cdot(\cos\lambda+i\sin\lambda)
=\operatorname{cotang}\tfrac12 P\left(\frac{1-e\cos P}{1+e\cos P}\right)^{\frac12 e}
\tag*{1)}
\end{equation*}
\begin{equation*}
x+iy=\int\frac{a(1-ee)\,\mathrm dP}{(1-ee\cos P^2)^{\frac32}},
\tag*{2)}
\end{equation*}
\noindent where \(a\) and \(a\surd(1-ee)\) are the two semiaxes of the ellipsoid
[and \(p=90^\circ-\Phi\) is set.

If we set
\[
\operatorname{cotang}\tfrac12 p\left(\frac{1-e\cos p}{1+e\cos p}\right)^{\frac12 e}
=\operatorname{cotang}\tfrac12 q
\]
\noindent and]
\[
\operatorname{cotang}\tfrac12 q(\cos\lambda+i\sin\lambda)=\operatorname{cotang}\tfrac12 Q,
\]
\noindent [as well as]
\[
Q=R+iS,
\]
\noindent\textsc{ix.}\hfill 25\hspace*{.12\linewidth}\par

% Printed 194; principal PDF 201; Google PDF 205.
\par\clearpage
\noindent\makebox[12mm][l]{\fontsize{13}{16}\selectfont 194}\makebox[\dimexpr\linewidth-24mm\relax][c]{\scriptsize\textbf{POSTHUMOUS PAPERS.\quad CONFORMAL MAPPING OF THE SPHEROID IN THE PLANE.}}\makebox[12mm][r]{}\par\medskip

\noindent [then we have]
\[
\begin{aligned}
\operatorname{tang}\tfrac12 q(\cos\lambda-i\sin\lambda)&=\operatorname{tang}\tfrac12(R+iS)\\
\operatorname{tang}\tfrac12 q(\cos\lambda+i\sin\lambda)&=\operatorname{tang}\tfrac12(R-iS).
\end{aligned}
\]
[From this it follows:
\[
\begin{alignedat}{2}
2\operatorname{tang}\tfrac12 q\cos\lambda&=\frac{\sin R}{\tfrac12(\cos R+\cos iS)},\qquad&
-2i\operatorname{tang}\tfrac12 q\sin\lambda&=\frac{\sin iS}{\tfrac12(\cos R+\cos iS)},\\
1-\operatorname{tang}\tfrac12 q^2&=\frac{\cos R}{\tfrac12(\cos R+\cos iS)},&
1+\operatorname{tang}\tfrac12 q^2&=\frac{\cos iS}{\tfrac12(\cos R+\cos iS)};
\end{alignedat}
\]
\noindent hence]
\[
\begin{aligned}
\operatorname{tang} R&=\frac{2\cos\lambda\operatorname{tang}\tfrac12 q}{1-\operatorname{tang}\tfrac12 q^2}
=\cos\lambda\operatorname{tang} q\\
\operatorname{tang} iS&=-\frac{2i\sin\lambda\operatorname{tang}\tfrac12 q}{1+\operatorname{tang}\tfrac12 q^2}
=-i\sin\lambda\sin q
\end{aligned}
\]
\noindent [or]
\[
S=\tfrac12\log\frac{1-\sin\lambda\sin q}{1+\sin\lambda\sin q}.
\]
\begin{center}\rule{.18\linewidth}{.4pt}\end{center}
\small
\setlength{\abovedisplayskip}{5pt}
\setlength{\belowdisplayskip}{5pt}
\setlength{\abovedisplayshortskip}{3pt}
\setlength{\belowdisplayshortskip}{3pt}

% Printed 195; principal PDF 202; Google PDF 206.
\par\clearpage

\thispagestyle{empty}
\vspace*{1.6\baselineskip}
\begin{center}\large\textbf{REMARKS.}\end{center}
{\small
A continuous account of the conformal mapping of the spheroid in the plane,
which served as the projection method for the Hanoverian arc measurement and land survey,
was not found among the posthumous papers. The preceding developments were compiled
with the aid of notes in two notebooks and a series of loose sheets which \textsc{Gauss}
may have used for calculations while deriving his formulas.

The notebook »Den astronomischen Wissenschaften gewidmet« contains the formulas for
calculating the geographic coordinates, meridian convergence and scale factor from
the plane rectangular coordinates. Articles [4], [9], [10], [11] were taken directly
from it, as were the articles on azimuth reduction [15], with the table for \(\log q\),
[16], [17], and those on the transformation of coordinates [18], [19], [20].
Articles [7], [13] and [14] were assembled from various passages in this notebook.

A second notebook, entitled »Kleine Aufsätze aus verschiedenen Theilen der
Mathematik«, supplied articles [5], [26], [27] and [31]. A loose sheet was also
used for [27].

Articles [1], [2], [6], [8], [21] and [22] were compiled from notes in the first
notebook and on loose sheets.

Articles [3], with its accompanying table, [12], [23], [29] and [32] were taken
from loose sheets, as was the second table for [15]. This table is reproduced here
in full because there is no other one based on a flattening different from
\(1:302{,}68\). This table is in a separate booklet. The headings of the individual
columns have been added here, as in the first table for [15]. \textsc{Gauss}
had written Corrige \(-26{,}264\) above the first column, which, as already
mentioned, corresponds to a decrease in the latitude of the origin by
\(0^{\prime\prime}{,}85\). The table for [28], of which only an extract has been
given above --- in the original it extends from \(\psi=5000000\,\mathrm m\) to
\(\psi=6500000\,\mathrm m\) --- and the table under [30] are on the same sheet.

Finally, articles [24], [25] and [28] were assembled from notes on various
sheets.

\medskip
A number of changes have been made in the individual articles.

\medskip
\noindent In [2]:\quad In place of the last formula for \(c\), \textsc{Gauss} has on a loose sheet
\[
c=\lambda\sin\varphi-
\frac{(1-ee\sin\varphi^2)^{\frac32}\operatorname{tang}\varphi}{6a^3(1-ee)^2}
\left\{1-ee+ee(3+ee)\cos\varphi^2+4e^4\cos\varphi^4\right\}y^3.
\]
\noindent »\quad [4]:\quad In the formula for \(c\), his expression for \(\log\operatorname{hyp} C'\) lacks the term
\[
-\tfrac18\frac{1-ee\sin\varphi^2}{aa}yy.
\]
\noindent »\quad [7]:\quad In place of the first formula the original reads
\[
N+i\log\mathrm{Conv.}=\log f'(\xi)-\log f'(x+iy),\qquad N=\log n,
\]
\noindent\hspace*{3.5em}which is probably just a slip of the pen.
\par}
\noindent\hfill \(25^{*}\)\hspace*{.12\linewidth}\par

% Printed 196; principal PDF 203; Google PDF 207.
\par\clearpage
\noindent\makebox[0pt][l]{\fontsize{13}{16}\selectfont 196}\hfill{\small\textbf{REMARKS.}}\hfill\null\par\medskip

{\small
\noindent In [8]:\quad The last formula is on a loose sheet, where it reads
\[
n=1-\frac{\gamma}{2\alpha}yy+
\frac{\alpha\alpha\varepsilon-6\alpha\beta\delta+6\beta\beta\gamma}{24\alpha^3}y^4.
\]
\noindent »\quad [9]:\quad In the formula for \(\log n\), \textsc{Gauss}'s coefficient of \(y^6\) lacks the term \(\tfrac1{720}A^{\mathrm{IV}}\).

\noindent »\quad [11]:\quad In place of the stated formulas for \(x\) and \(y\), the original reads:
\[
\begin{aligned}
x&=\xi-\tfrac12\alpha\beta\lambda\lambda
+\tfrac1{24}(\alpha\beta^3+4\alpha\alpha\beta\gamma+\alpha^3\delta)\lambda^4\\
y&=\alpha\lambda-\tfrac16(\alpha\alpha\gamma+\alpha\beta\beta)\lambda^3.
\end{aligned}
\]
\noindent »\quad [12]:\quad The factor of \(y^6\) in the formula for \(\xi\) is given as follows on a slip of paper that lay on \textsc{Gauss}'s desk with other slips at his death:
\[
\tfrac1{45}\alpha_1^5+\tfrac5{144}\alpha_1^3\beta_1
-\tfrac1{144}\alpha_1\alpha_1\gamma_1+\tfrac1{48}\alpha_1\beta_1\beta_1
+\tfrac1{72}\beta_1\gamma_1-\tfrac1{144}\alpha_1\delta_1+\tfrac1{720}\varepsilon_1.
\]
\noindent »\quad [13]:\quad In the expression for \(\mathrm g^{\mathrm V}(\mathrm F)\),
and correspondingly in the first formula for \(y\), in the factor of \(\lambda^5\)
on p.\ 157, the value in parentheses in the original reads:
\[
(1-20cc+(24-57\delta)c^4+\ldots),
\]
\noindent whereas it should read as stated above:
\[
(1-20cc+(24-58\delta)c^4+\ldots).
\]
\noindent »\quad [15]:\quad In place of \(\log\mathfrak{B}=7{,}317\,8248-30\),
\textsc{Gauss} has \(\log\mathfrak{B}=7{,}317\,8148-30\). Consequently, in the
second table of Art.\ [15], pp.\ 161 and 162, the values for \(\log\beta\)
are to be increased by one unit in the last place.

\noindent »\quad [16]:\quad The coefficient of \(R^3\) in the formula for \(-n^*\sin\varphi\),
p.\ 164, reads in the original
\[
\tfrac14c'+\frac{ab'}{24N}-\frac{a'b''}{24N}
+\frac{a^{\prime3}}{48NN}+\frac{aaa'}{24NN}.
\]
\noindent In the formula for \(\log n^*\) at the top of p.\ 165, \textsc{Gauss}
writes the terms of second degree as
\[
-Cxx-2Dxy-Eyy;
\]
\noindent but the expansion that follows disagrees with this.

\noindent »\quad [17]:\quad For \(r\), p.\ 168, the original gives
\[
r=R+\tfrac23ttR^3+(\tfrac32tu-\tfrac12t^4)R^4\ldots;
\]
\noindent so the resulting differential equation for \(\psi\) is
\[
\begin{aligned}
&\left(1+(2tt+t't')RR+(6tu-2t^4+2t'u')R^3\ldots\right)
 \frac{\partial\psi}{\partial R}\\
&\qquad+\left(-\tfrac23tt'R-(\tfrac12tu'+\tfrac32t'u+2t't^3)RR\ldots\right)
 \frac{\partial\psi}{\partial\theta}=0,
\end{aligned}
\]
\noindent from which we then obtain:
\[
\psi=\theta+\tfrac13tt'RR+\left((\tfrac12u+\tfrac23t^3)t'+\tfrac16tu'\right)R^3\ldots.
\]
\noindent »\quad [20]:\quad In the formula for \(x+\Omega\), \textsc{Gauss} has as the factor of \(\omega\)
\[
1-\tfrac12(1+tt)yy+(\tfrac5{24}-\tfrac{13}{24}tt-\tfrac{15}8t^4)y^4\ldots.
\]
\noindent »\quad [23]:\quad In the original the coefficient of \(y^9\) in the formula for \(x\), p.\ 175, reads
\[
\frac{18621\delta^4+48826\delta^5+48160\delta^6+21288\delta^7+3548\delta^8}{2835},
\]
\noindent and correspondingly the coefficient of \(y^{10}\) in the formula for \(xx\) reads:
\[
\frac{316305\delta^4+831340\delta^5+822217\delta^6+363738\delta^7+60623\delta^8}{14175};
\]
\par}

% Printed 197; principal PDF 204; Google PDF 208.
\par\clearpage
\noindent\null\hfill{\small\textbf{CONFORMAL MAPPING OF THE SPHEROID IN THE PLANE.}}\hfill\makebox[0pt][r]{\fontsize{13}{16}\selectfont 197}\par\medskip

{\small
\noindent In the next formula, for \(1-xx\), the original gives the following as the coefficient of \(y^8\):
\[
\frac{7029\delta^4+16182\delta^5+21574\delta^6+5798\delta^7+892\delta^8}{315}.
\]
\noindent In [24]:\quad The factor of \(\sin\varphi^5\cos\varphi\) in the second formula
for \(\psi\) at the top of p.\ 177 reads in the original:
\[
\tfrac{241}{480}e^6-\tfrac{3777}{30420}e^8.
\]
\noindent As the expression for \(\sin\varphi\cos\varphi\), p.\ 177, \textsc{Gauss} gives:
\[
\sin\omega\cos\omega\left(1+ee+e^4+e^6-\frac{12ee+25e^4+39e^6}{6}yy
+\frac{65e^4+258e^6}{15}y^4-\frac{146}{15}e^6y^6\right).
\]
\noindent The formula for \(\psi\) expressed in terms of \(\omega\), at the bottom of p.\ 177, reads in the original:
\[
\begin{aligned}
\psi=\omega&+\sin2\omega\left(\frac18ee-\frac{13}{192}e^4+\frac{173}{1024}e^6
+\frac{29029}{368640}e^8\ldots\right)\\
&-\sin4\omega\left(\frac{47}{768}e^4-\frac{293}{3840}e^6+\frac{1220619}{884736}e^8\ldots\right)\\
&+\sin6\omega\left(\frac{979}{15360}e^6-\frac{4124089}{10321920}e^8\ldots\right)\\
&-\sin8\omega\left(\frac{12895751}{123863040}e^8\ldots\right).
\end{aligned}
\]
\noindent In \textsc{Gauss}'s text the factor of \(-\sin2\varphi\) in the last formula,
which gives \(\omega\) as a function of \(\varphi\), p.\ 178, reads
\[
\tfrac12ee+\tfrac5{24}e^4+\tfrac3{32}e^6+\tfrac{281}{11520}e^8.
\]
\noindent »\quad [26]:\quad In place of the first two lines of the expression
for \(\log\cos\varphi\), the original reads:
\[
\begin{aligned}
\log\cos\varphi&=\log\cos\psi-6ff-63f^4-\frac{2192}3f^6-\frac{17559}2f^8\ldots\\
&\qquad+\left(6ff+42f^4+\frac{946}3f^6+2200f^8\ldots\right)\cos2\psi.
\end{aligned}
\]
\noindent »\quad [28]:\quad In the formula expressing \(\psi\) as a function
of \(\omega\), p.\ 183, \textsc{Gauss} has
\[
4{,}800\,0528\sin4\omega
\]
\[
\text{in place of }4{,}700\,0528\sin4\omega,
\]
\noindent and correspondingly, in the calculation in [29], p.\ 188,
\(\log B=0{,}681\,2460\) in place of \(\log B=0{,}672\,1028\). His value of \(y\),
p.\ 188, is thus \(0{,}007\,\mathrm m\) smaller than the value stated above.

\noindent The value used in the developments under [28]
\[
\log\frac{2\,.\,10^7}{\pi}ff=3{,}418\,9411.803.
\]
\noindent which also underlies the table on p.\ 184 for passing from \(\psi\) to \(\varphi\)
and \(\omega\), probably arose from a slip of the pen by \textsc{Gauss}. On a loose
sheet the value of \(\log\frac{2\,.\,10^7}{\pi}ff\) just given is accompanied by
\(n=302{,}78\). That assignment is incorrect. The value corresponding to this
flattening is
\[
\log\frac{2\,.\,10^7}{\pi}ff=3{,}418\,9451.803.
\]
\par}

% Printed 198; principal PDF 205; Google PDF 209.
\par\clearpage
\noindent\makebox[0pt][l]{\fontsize{13}{16}\selectfont 198}\hfill{\small\textbf{REMARKS.}}\hfill\null\par\medskip

{\small
\noindent In [31]:\quad In the second formula for \(\log\operatorname{tang}\tfrac12 p\), p.\ 191, the original reads:
\[
\begin{array}{cccc@{\quad}ccc@{\quad}rcl}
\text{in the}&\text{coefficient}&\text{of}&-\cos q&
 \text{in}&\text{place}&\text{of}&\tfrac{11}{48}a^3&:&\tfrac{11}{24}a^3\\
\text{»}&\text{»}&\text{»}&-\cos3q&
 \text{»}&\text{»}&\text{»}&\tfrac3{32}a^3&:&\tfrac3{16}a^3\\
\text{»}&\text{»}&\text{»}&-\cos5q&
 \text{»}&\text{»}&\text{»}&\tfrac1{96}a^3&:&\tfrac1{48}a^3.
\end{array}
\]
\noindent Consequently \textsc{Gauss}'s text reads:
\[
\begin{aligned}
\log\operatorname{tang}(45^\circ+\tfrac12\omega)
 &=\log\operatorname{tang}(45^\circ+\tfrac12\psi)
 -\left(4ff+4f^4+\frac{40}3f^6-22f^8\ldots\right)\sin\psi\\
 &\quad-\left(\frac83f^4+\frac{34}3f^6+59f^8\ldots\right)\sin3\psi
 -\left(\frac{98}{15}f^6+\frac{247}3f^8\ldots\right)\sin5\psi\\
 &\quad-\left(\frac{400}{21}f^8\ldots\right)\sin7\psi.
\end{aligned}
\]
\noindent As numerical values of the coefficients of this series, \textsc{Gauss} gives:
\[
\begin{array}{r@{\quad}l@{\quad}l@{\quad}l}
10494{,}784\,2962\,\mathrm m;&2{,}888\,891\,3195\,\mathrm m;&
0{,}002\,927\,170\,360\,\mathrm m;&0{,}000\,003\,499\,170\,05\,\mathrm m;\\
340^{\prime\prime}{,}031\,0112;&0^{\prime\prime}{,}093\,600\,078\,75;&
0^{\prime\prime}{,}000\,094\,840\,32;&0^{\prime\prime}{,}000\,000\,113\,373\,11;
\end{array}
\]
\noindent and as final formulas:
\[
\begin{aligned}
\lambda={}&\operatorname{arc}\operatorname{tang}\frac{\sin iy}{i\cos x}\\
&-340^{\prime\prime}{,}031\,0112\cos x\frac{\sin iy}{i}
 -0^{\prime\prime}{,}093\,6001\cos3x\frac{\sin3iy}{i}\\
&-0^{\prime\prime}{,}000\,0948\cos5x\frac{\sin5iy}{i}
 -0^{\prime\prime}{,}000\,0001\cos7x\frac{\sin7iy}{i},
\end{aligned}
\]
\[
\begin{aligned}
\log\operatorname{tang}(45^\circ+\tfrac12\Omega)
={}&\log\surd\frac{\cos iy+\sin x}{\cos iy-\sin x}\\
&-0{,}000\,715\,9403\sin x\cos iy
 -0{,}000\,000\,1971\sin3x\cos3iy\\
&-0{,}000\,000\,0002\sin5x\cos5iy.
\end{aligned}
\]
\noindent In [32]:\quad The expression for \(S\), p.\ 194, lacks the factor \(\tfrac12\) in the original.
\begin{center}\rule{.24\linewidth}{.4pt}\end{center}

The following further remarks concern the formulas themselves.

The basic formulas of \textsc{Gauss}'s development, Art.\ [1], are obtained by
mapping the plane conformally onto the spheroid. The problem is thus the inverse
of the mapping treated in Art.\ 12 of the »Allgemeinen Auflösung der Aufgabe:
Die Theile einer gegebnen Fläche auf einer andern gegebnen Fläche so abzubilden,
dass die Abbildung dem Abgebildeten in den kleinsten Theilen ähnlich wird«,
Vol.\ IV, p.\ 205 ff. Here there is the additional condition that segments on
the abscissa axis in the plane must equal the corresponding arcs of a particular
meridian, the principal meridian. If \(x,y\) are the rectangular coordinates of a point
in the plane, and \(\Phi,\lambda\) the geographic coordinates of the corresponding
point on the spheroid, then following Art.\ 12 we have:
\[
\omega=\mathrm dx^2+\mathrm dy^2,
\]
\noindent and the differential equation \(\omega=0\) gives \(x\pm iy=\mathrm{const.}\)

Further,
\[
\begin{aligned}
\Omega&=\frac{aa(1-ee)^2}{(1-ee\sin\Phi^2)^3}\,\mathrm d\Phi^2
 +\frac{aa\cos\Phi^2}{1-ee\sin\Phi^2}\,\mathrm d\lambda^2\\
&=\frac{aa\cos\Phi^2}{1-ee\sin\Phi^2}
\left\{\left(\frac{1-ee}{(1-ee\sin\Phi^2)\cos\Phi}\right)^2
 \mathrm d\Phi^2+\mathrm d\lambda^2\right\};
\end{aligned}
\]
\noindent so the differential equation \(\Omega=0\) gives
\[
\frac{1-ee}{(1-ee\sin\Phi^2)\cos\Phi}\,\mathrm d\Phi\pm i\,\mathrm d\lambda=0,
\]
\par}

% Printed 199; principal PDF 206; Google PDF 210.
\par\clearpage
\noindent\null\hfill{\small\textbf{CONFORMAL MAPPING OF THE SPHEROID IN THE PLANE.}}\hfill\makebox[0pt][r]{\fontsize{13}{16}\selectfont 199}\par\medskip

{\small
\noindent or, if we set
\begin{equation*}
\frac{(1-ee\sin\Phi^2)\cos\Phi}{1-ee}=\Theta(\Phi)
\tag*{1)}
\end{equation*}
\noindent then we have
\[
\frac{\mathrm d\Phi}{\Theta(\Phi)}\pm i\,\mathrm d\lambda=0,
\]
\noindent whose integral is
\[
\int\frac{\mathrm d\Phi}{\Theta(\Phi)}\pm i\lambda=\mathrm{const.},
\]
\noindent or, if we set
\begin{equation*}
\mathrm F(\Phi)=\int\frac{\mathrm d\Phi}{\Theta(\Phi)}
=\log\operatorname{hyp}\operatorname{tang}(45^\circ+\tfrac12\Phi)
+\tfrac12 e\log\operatorname{hyp}\frac{1-e\sin\Phi}{1+e\sin\Phi}
\tag*{2)}
\end{equation*}
\noindent the equation
\[
\mathrm F(\Phi)\pm i\lambda=\mathrm{const.}
\]
\noindent holds.

The latitude \(\Phi\) is to correspond to the meridian arc \(\xi\), measured from the
equator, so that the equation
\begin{equation*}
\mathrm F(\Phi)=f(\xi)=\int\frac{\mathrm d\Phi}{\Theta(\Phi)}
\tag*{3)}
\end{equation*}
\noindent also holds.

If \(\mathfrak F\) denotes an arbitrary function, the plane is thus mapped
conformally onto the ellipsoid by the equation
\[
f(\xi)+i\lambda=\mathfrak F(x+iy).
\]
Now each segment on the abscissa axis is to equal the corresponding meridian arc.
Thus, for \(\lambda=0\), we must have \(\xi=x\) and \(y=0\), from which it follows that
\[
f(x)=\mathfrak F(x).
\]
\noindent Thus the mapping equation is:
\begin{equation*}
f(\xi)+i\lambda=\mathrm F(\Phi)+i\lambda=f(x+iy).
\tag*{4)}
\end{equation*}
\begin{center}\rule{.24\linewidth}{.4pt}\end{center}

The formulas for the scale factor, Arts.\ [6] to [10], may be derived as follows.

Since
\[
\begin{gathered}
\mathrm df(\xi)=\frac{\mathrm d\Phi}{\Theta(\Phi)}
=\frac{(1-ee)\,\mathrm d\Phi}{\cos\Phi(1-ee\sin\Phi^2)}\\
f'(\xi)=\frac{1-ee}{\cos\Phi(1-ee\sin\Phi^2)}
\,.\,\frac{(1-ee\sin\Phi^2)^{\frac32}}{a(1-ee)}
=\frac{\surd(1-ee\sin\Phi^2)}{a\cos\Phi},
\end{gathered}
\]
\noindent where \(\xi\) increases with \(\Phi\) and is thus taken to be positive northward, we obtain
\[
\begin{aligned}
\frac1{nn}=\frac{\Omega}{\omega}
 &=\left(\frac1{f'(\xi)}\right)^2
 \frac{(\mathrm df(\xi))^2+\mathrm d\lambda^2}{\mathrm dx^2+\mathrm dy^2}
 =\left(\frac1{f'(\xi)}\right)^2
 \frac{\mathrm d(f(\xi)+i\lambda)}{\mathrm d(x+iy)}
 \,.\,\frac{\mathrm d(f(\xi)-i\lambda)}{\mathrm d(x-iy)}\\
 &=\left(\frac1{f'(\xi)}\right)^2
 \left\{\left(\frac{\partial\lambda}{\partial x}\right)^2+
 \left(\frac{\partial\lambda}{\partial y}\right)^2\right\}.
\end{aligned}
\]
Introducing the meridian convergence \(c\), we first have
\(\operatorname{tang}c=-\dfrac{\mathrm dy}{\mathrm dx}\), where
\(\dfrac{\mathrm dy}{\mathrm dx}\) is obtained from the equation
\(\lambda=\mathrm{const.}\). Hence
\par}

% Printed 200; principal PDF 207; Google PDF 211.
\par\clearpage
\noindent\makebox[0pt][l]{\fontsize{13}{16}\selectfont 200}\hfill{\small\textbf{REMARKS.}}\hfill\null\par\medskip

{\small
\[
\frac1{\cos c^2}=\frac{\left(\dfrac{\partial\lambda}{\partial x}\right)^2
 +\left(\dfrac{\partial\lambda}{\partial y}\right)^2}
 {\left(\dfrac{\partial\lambda}{\partial y}\right)^2},\qquad
\frac1{\sin c^2}=\frac{\left(\dfrac{\partial\lambda}{\partial x}\right)^2
 +\left(\dfrac{\partial\lambda}{\partial y}\right)^2}
 {\left(\dfrac{\partial\lambda}{\partial x}\right)^2},
\]
\noindent and therefore
\begin{equation*}
\GaussOldenburgLeftBrace{\begin{aligned}
\frac{\partial\lambda}{\partial y}
 &=\phantom{-}\frac{\partial f(\xi)}{\partial x}=\frac1n f'(\xi)\cos c=u\\
\frac{\partial\lambda}{\partial x}
 &=-\frac{\partial f(\xi)}{\partial y}=\frac1n f'(\xi)\sin c=v.
\end{aligned}}
\tag*{1)}
\end{equation*}
These formulas follow more simply by observing that, in the plane (assuming
that \(x\) is positive northward and \(y\) increases with \(\lambda\)), for the
image of an element of the parallel the relation
\begin{equation*}
\frac{n\,.\,a\cos\Phi}{\surd(1-ee\sin\Phi^2)}\,\mathrm d\lambda
=\frac n{f'(\xi)}\,\mathrm d\lambda=\mathrm dx\sin c+\mathrm dy\cos c
\tag*{2)}
\end{equation*}
\noindent holds.

From 1) we obtain:
\begin{equation*}
\frac1{nn}=\left(\frac u{f'(\xi)}\right)^2+\left(\frac v{f'(\xi)}\right)^2.
\tag*{3)}
\end{equation*}
Using this formula, \textsc{Gauss} calculated the scale factor in [6].

The same factor could have been obtained somewhat more easily using the formula that follows from 3),
\begin{equation*}
\frac1{nn}=\left(\frac{\partial\xi}{\partial x}\right)^2+
\left(\frac{\partial\xi}{\partial y}\right)^2
\tag*{4)}
\end{equation*}
\noindent This formula follows directly by writing, for the image of an element
of the meridian,
\[
\frac{n\,.\,a(1-ee)}{\surd(1-ee\sin\Phi^2)^3}\,\mathrm d\Phi
=n\,\mathrm d\xi=\mathrm dx\cos c-\mathrm dy\sin c,
\]
\noindent from which it follows that
\[
\frac{\partial\xi}{\partial x}=\frac1n\cos c,\qquad
\frac{\partial\xi}{\partial y}=-\frac1n\sin c.
\]
By 1) we further have
\[
\frac1{f'(\xi)}\left\{\frac{\partial f(\xi)}{\partial x}
+i\frac{\partial\lambda}{\partial x}\right\}
=\frac1n(\cos c+i\sin c)
\]
\noindent or
\[
\frac{f'(x+iy)}{f'(\xi)}=\frac1n e^{ic},
\]
\noindent or, if we set
\[
\log n=N
\]
\noindent then
\begin{equation*}
N-ic=\log f'(\xi)-\log f'(x+iy).
\tag*{5)}
\end{equation*}
The resulting value
\[
N=\log f'(\xi)-\text{Real part of }\log f'(x+iy)
\]
\noindent is also obtained immediately from the first formula given on the preceding page.
From that formula we have
\[
n=\frac{f'(\xi)}{\surd\{f'(x+iy)f'(x-iy)\}}.
\]
Differentiating equation 5) twice with respect to \(x\) and \(y\), we obtain
the following by 4),
\par}

% Printed 201; principal PDF 208; Google PDF 212.
\par\clearpage
\noindent\null\hfill{\small\textbf{CONFORMAL MAPPING OF THE SPHEROID IN THE PLANE.}}\hfill\makebox[0pt][r]{\fontsize{13}{16}\selectfont 201}\par\medskip

{\small
\noindent taking into account that
\[
\frac{\partial\partial f(\xi)}{\partial x^2}
+\frac{\partial\partial f(\xi)}{\partial y^2}=0
\]
\noindent holds:
\begin{equation*}
\frac{\partial\partial c}{\partial x^2}
+\frac{\partial\partial c}{\partial y^2}=0
\tag*{6)}
\end{equation*}
\noindent and
\begin{equation*}
\frac{\partial\partial N}{\partial x^2}
+\frac{\partial\partial N}{\partial y^2}
=\frac1{nn}\frac{f'''(\xi)f'(\xi)-2(f''(\xi))^2}{(f'(\xi))^2}
=-\frac1{nn}\frac{\theta''(\xi)}{\theta(\xi)},
\tag*{7)}
\end{equation*}
\noindent where the relation
\[
\theta(\xi)=\frac1{f'(\xi)}=\frac{a\cos\Phi}{\surd(1-ee\sin\Phi^2)}
\]
\noindent holds.

Since \(\theta'(\xi)=-\sin\Phi\) and
\(\theta''(\xi)=-\cos\Phi\dfrac{(1-ee\sin\Phi^2)^{\frac32}}{a(1-ee)}\) by the above, and since
\[
r=\frac{a(1-ee)}{\surd(1-ee\sin\Phi^2)^3},\qquad
r'=\frac{a}{\surd(1-ee\sin\Phi^2)},
\]
\noindent we obtain
\begin{equation*}
\frac{\partial\partial N}{\partial x^2}
+\frac{\partial\partial N}{\partial y^2}=\frac1{rr'nn}.
\tag*{8)}
\end{equation*}
\noindent Cf.\ Vol.\ VIII, p.\ 385.
\begin{center}\rule{.24\linewidth}{.4pt}\end{center}

The following should also be noted about the azimuth reduction given in Arts.\ [15]
and [16] (cf.\ the letter to \textsc{Schumacher} of 25 June 1831).
By the plane azimuth of the point \(x_2,y_2\) at the point \(x_1,y_1\), \textsc{Gauss}
means the angle that the straight line joining the two points makes with the line
through \(x_1,y_1\) parallel to the \(x\)-axis. The azimuth on the spheroid, on the other
hand, is the angle that the geodesic at the point corresponding to \(x_1,y_1\) makes with
the curve whose image in the plane is a line parallel to the abscissa axis.
If \(x\) is positive southward, and if at the point \((x_1y_1)\)
\(T\) denotes the astronomical azimuth, \(t\) the azimuth on the spheroid, \(\theta\) the
plane azimuth, \(c\) the meridian convergence, and \(\psi\) the angle that the image
of the geodesic makes with the straight line through \((x_1y_1)\) and \((x_2y_2)\), then
\[
T=t-c,\qquad \theta=t-\psi.
\]
The formulas listed under [15] for reducing the azimuth on the spheroid to the
plane azimuth follow from the equation for \(\psi\) derived in Art.\ [16], p.\ 166,
when its terms of second order are neglected. We then have
\[
\psi=\tfrac12 B_{\frac13}\,.\,(x_2-x_1)-\tfrac12 A_{\frac13}\,.\,(y_2-y_1).
\]
The coordinates \(x_1,y_1\) correspond to the starting point and the coordinates \(x_2,y_2\)
to the endpoint of the geodesic. To obtain the coefficients \(A_{\frac13}\) and
\(B_{\frac13}\), \(\log n\) is to be expressed as a function of the coordinates.
Let \(x,y\) be any point to which a point at latitude \(\varphi\) corresponds
on the ellipsoid. If we then set
\begin{equation*}
\tfrac12\frac{(1-ee\sin\varphi^2)^2}{aa(1-ee)}=\alpha
\tag*{1)}
\end{equation*}
\par}
\noindent\textsc{ix.}\hfill 26\hspace*{.12\linewidth}\par
\endgroup
